% TOP_PROBLEM: {"problemId": "problem.non-triviality-of-the-three-dimensional-ising-model", "canonicalTitle": "Non-Triviality of the Three-Dimensional Ising Model", "releaseRank": 73, "primaryDomain": "mathematical_physics", "primaryDomainLabel": "Mathematical physics", "displayStatus": "open", "releaseStatus": "open"}
% !TeX program = lualatex
% Definition and sources only; this file is not an LLM solution attempt.
\documentclass[11pt]{article}
\usepackage[margin=25mm]{geometry}
\usepackage{fontspec}
\setmainfont{Latin Modern Roman}
\usepackage{amsmath,amssymb,mathtools}
\usepackage{unicode-math}
\setmathfont{Latin Modern Math}
\usepackage{xurl,hyperref}
\hypersetup{colorlinks=true,urlcolor=blue}
\setcounter{secnumdepth}{0}
\setlength{\parindent}{0pt}
\setlength{\parskip}{5pt}
\setlength{\emergencystretch}{3em}
\begin{document}
\section*{\texorpdfstring{Non-Triviality of the Three-Dimensional Ising Model}{Non-Triviality of the Three-Dimensional Ising Model}}\label{73}
\subsection{Definitions and mathematical statement}
At the critical point of the nearest-neighbor three-dimensional Ising model, consider suitably normalized spin fields under vanishing lattice spacing. A centered Gaussian field has all odd moments zero and its even moments equal sums over pairings of two-point correlations. Nontriviality in this record means a scaling-limit field violating this rule; a possible witness is a nonzero truncated four-point correlation
\[
 \langle\Phi_1\Phi_2\Phi_3\Phi_4\rangle
 -\langle\Phi_1\Phi_2\rangle\langle\Phi_3\Phi_4\rangle
 -\langle\Phi_1\Phi_3\rangle\langle\Phi_2\Phi_4\rangle
 -\langle\Phi_1\Phi_4\rangle\langle\Phi_2\Phi_3\rangle.
\]
Conformal invariance is not part of this narrower target. The normalization and mode of convergence must be specified in a rigorous formulation of the scaling limit.
\par\smallskip\textit{Formulation and concept references:} \href{https://arxiv.org/html/1707.00520}{[S1]}.

\subsection{Short English statement}
Can one prove that the critical three-dimensional Ising scaling limit is genuinely interacting rather than Gaussian?
\subsection{Sources}
\begin{itemize}
\item[S1]Lectures on the Ising and Potts models on the hypercubic lattice (Hugo Duminil-Copin).
\url{https://arxiv.org/html/1707.00520}.
\textit{Formulation source.}
\item[S2]Triviality of the scaling limits of critical Ising and phi\textasciicircum{}4 models with effective dimension at least four.
\url{https://arxiv.org/abs/2309.05797}.
\textit{Further reference; not a proof endorsement.}
\item[S3]Lectures on the Ising and Potts models on the hypercubic lattice.
\url{https://www.ihes.fr/~/duminil/publi/2017PIMS.pdf}.
\textit{Further reference; not a proof endorsement.}
\item[S4]Certified three-dimensional Ising critical exponents from harmonized bootstrap series and interval Monte Carlo.
\url{https://link.springer.com/article/10.1140/epjb/s10051-025-01109-8}.
\textit{Further reference; not a proof endorsement.}
\item[S5]Certified Three Dimensional Ising Critical Exponents: author release notes.
\url{https://www.researchgate.net/publication/397730833_Certified_Three_Dimensional_Ising_Critical_Exponents}.
\textit{Further reference; not a proof endorsement.}
\end{itemize}
\end{document}
